% Generated by scripts/build_reports.py; edit reviewed Markdown inputs.
% Compile with XeLaTeX.
% Options for packages loaded elsewhere
\PassOptionsToPackage{unicode}{hyperref}
\PassOptionsToPackage{hyphens}{url}
\PassOptionsToPackage{dvipsnames,svgnames,x11names}{xcolor}
\documentclass[
  11pt,
]{article}
\usepackage{xcolor}
\usepackage[margin=1in]{geometry}
\usepackage{amsmath,amssymb}
\setcounter{secnumdepth}{-\maxdimen} % remove section numbering
\usepackage{iftex}
\ifPDFTeX
  \usepackage[T1]{fontenc}
  \usepackage[utf8]{inputenc}
  \usepackage{textcomp} % provide euro and other symbols
\else % if luatex or xetex
  \usepackage{unicode-math} % this also loads fontspec
  \defaultfontfeatures{Scale=MatchLowercase}
  \defaultfontfeatures[\rmfamily]{Ligatures=TeX,Scale=1}
\fi
\usepackage{lmodern}
\ifPDFTeX\else
  % xetex/luatex font selection
\fi
% Use upquote if available, for straight quotes in verbatim environments
\IfFileExists{upquote.sty}{\usepackage{upquote}}{}
\IfFileExists{microtype.sty}{% use microtype if available
  \usepackage[]{microtype}
  \UseMicrotypeSet[protrusion]{basicmath} % disable protrusion for tt fonts
}{}
\makeatletter
\@ifundefined{KOMAClassName}{% if non-KOMA class
  \IfFileExists{parskip.sty}{%
    \usepackage{parskip}
  }{% else
    \setlength{\parindent}{0pt}
    \setlength{\parskip}{6pt plus 2pt minus 1pt}}
}{% if KOMA class
  \KOMAoptions{parskip=half}}
\makeatother
\usepackage{longtable,booktabs,array}
\usepackage{caption}
\captionsetup[table]{skip=6pt}
\newcounter{none} % for unnumbered tables
\usepackage{calc} % for calculating minipage widths
% Correct order of tables after \paragraph or \subparagraph
\usepackage{etoolbox}
\makeatletter
\patchcmd\longtable{\par}{\if@noskipsec\mbox{}\fi\par}{}{}
\makeatother
% Allow footnotes in longtable head/foot
\IfFileExists{footnotehyper.sty}{\usepackage{footnotehyper}}{\usepackage{footnote}}
\makesavenoteenv{longtable}
\setlength{\emergencystretch}{3em} % prevent overfull lines
\providecommand{\tightlist}{%
  \setlength{\itemsep}{0pt}\setlength{\parskip}{0pt}}
\usepackage{amsmath,amssymb}
\usepackage{xurl}
\usepackage{fancyhdr}
\pagestyle{fancy}
\fancyhf{}
\fancyhead[L]{\small Pure Thought | Revised research specifications}
\fancyfoot[L]{\small 2026-09-17}
\fancyfoot[R]{\small\thepage}
\setlength{\headheight}{15pt}
\setlength{\emergencystretch}{3em}
\allowdisplaybreaks
\usepackage{bookmark}
\IfFileExists{xurl.sty}{\usepackage{xurl}}{} % add URL line breaks if available
\urlstyle{same}
% fallback for those not using the hyperref driver hyperxmp:
\makeatletter
\@ifundefined{xmpquote}{\newcommand{\xmpquote}[1]{#1}}{}
\makeatother
\hypersetup{
  pdftitle={S09. Stoquasticity under a bounded class of basis changes},
  pdfauthor={Pure Thought research portfolio},
  colorlinks=true,
  linkcolor={blue},
  filecolor={Maroon},
  citecolor={Blue},
  urlcolor={blue},
  pdfcreator={LaTeX via pandoc}}

\title{S09. Stoquasticity under a bounded class of basis changes}
\author{Pure Thought research portfolio}
\date{Revised 17 September 2026}

\begin{document}
\maketitle
\thispagestyle{empty}
\clearpage

\section{S09. Stoquasticity under a bounded class of basis changes}

\textbf{Source:} synthesis challenge 25;
\href{https://github.com/rick-stevens-ai/Pure-Thought/blob/main/NOTES/30-Most-Compelling-Pure-Thought-AI-Challenges.md}{original
synthesis}. This topic is distinct from, or adds a distinct objective
to, the canonical PRD catalog.
\href{https://github.com/rick-stevens-ai/Pure-Thought/blob/main/reviews/SYNTHESIS-CROSSWALK.md}{Crosswalk}

\subsection{Plain-language guide}

\subsubsection{The problem in everyday terms}

Quantum calculations can involve large positive and negative
contributions that almost cancel, making some sampling methods
inefficient. Sometimes a change in the mathematical basis removes
troublesome signs. This problem asks whether that is possible within a
finite, explicitly allowed set of local basis changes, so that both
success and failure can be checked completely.

\subsubsection{Key terms}

\begin{itemize}
\item
  \textbf{Basis} --- The reference states used to write a quantum
  operator as a matrix.
\item
  \textbf{Hamiltonian} --- The matrix describing the system's energy and
  couplings.
\item
  \textbf{Off-diagonal entry} --- A matrix element linking two different
  basis states.
\item
  \textbf{Stoquastic} --- Having real, nonpositive off-diagonal
  Hamiltonian entries in the specified basis.
\item
  \textbf{Sign problem} --- A sampling difficulty caused by signed or
  complex weights and cancellation in a chosen representation.
\item
  \textbf{Pauli operator} --- One of the standard single-qubit matrices
  used to express spin interactions.
\item
  \textbf{Clifford transformation} --- A basis transformation that maps
  Pauli operators to Pauli operators under conjugation.
\item
  \textbf{Product transformation} --- A change of basis applied
  independently to the individual qubits.
\item
  \textbf{Satisfiability certificate} --- A checkable assignment meeting
  encoded constraints, or an appropriate proof that none exists.
\end{itemize}

\subsubsection{Why this matters}

Knowing when signs can be removed helps separate difficulties caused by
representation from those that survive a proposed remedy. A bounded
search also creates small, exact benchmarks for algorithms intended to
find useful basis changes.

\subsubsection{What progress would mean}

A successful result would supply and verify a transformation, or exclude
every transformation in the allowed finite class. The complete
Hamiltonian must be checked because contributions can cancel.
Stoquasticity alone does not guarantee fast sampling, and failure for
product Clifford changes does not exclude all possible bases.

\subsection{Revised core question}

For a finite rational Pauli Hamiltonian, is there an allowed product
Clifford transformation making all off-diagonal entries real and
nonpositive?

\subsection{Precise scope, target and outputs}

Fix n qubits, rational coefficients, and allowed
U=\ensuremath{\otimes}\_iU\_i with U\_i from the finite single-qubit
Clifford group modulo phase. Define stoquasticity in the computational
basis for the full transformed H, not automatically term by term.
Exhaustively enumerate or use an exactly equivalent constraint encoding.
Export U and an exact matrix/sign check if successful; otherwise export
complete enumeration or a checked UNSAT proof. Continuous local
unitaries require a semialgebraic formulation and are a distinct
extension.

\subsection{Detailed critique of the synthesis framing}

The original omits which local basis changes are allowed, so
``unavoidable sign problem'' has no definite quantifier. A SAT search
over finitely many transformations cannot exclude all continuous local
rotations. Requiring every term to be stoquastic can also be stronger
than requiring their sum to be stoquastic because entries can cancel.
Finally, stoquasticity does not guarantee efficient sampling or mixing.
The rewrite makes a finite, falsifiable basis-cure problem and limits
every no-go conclusion to its transformation class.

\subsection{Validation and milestones}

Baseline: exhaustive small-n decisions with explicit transformations or
coverage. Strong: verified SAT encoding checked against brute force.
Extend to a continuous rotation class only with exact equations,
inequalities, and global proof machinery.

Use the shared
\href{https://github.com/rick-stevens-ai/Pure-Thought/blob/main/reviews/EVIDENCE-STANDARD.md}{evidence
standard}. Unsuccessful search is unresolved unless a complete exclusion
proof is supplied. These are proposed baseline specifications, not
completed scientific results.

\href{https://pmc.ncbi.nlm.nih.gov/articles/PMC6450938/}{Marvian, Lidar
and Hen, On the computational complexity of curing non-stoquastic
Hamiltonians} supports the need to specify the allowed basis class.

\subsection{Displayed formulation and claim boundary}

\href{https://github.com/rick-stevens-ai/Pure-Thought/blob/main/supplemental/S09.md}{Supplemental
specification}.

\[U=\bigotimes_{i=1}^nU_i,\quad U_i\in\mathrm{Clifford}_1/\mathrm{phase},
\qquad H'=UHU^\dagger,\]

\[\operatorname{Im}H'_{xy}=0,\qquad\operatorname{Re}H'_{xy}\le0\quad(x\ne y).\]

Apply these conditions to the full transformed matrix, allowing
cancellations between terms. Exhaustive coverage or an equivalent
checked UNSAT proof excludes only this finite transformation class.
Stoquasticity alone does not guarantee efficient mixing or sampling.

\section{Evidence and execution standard}

This standard applies to the revised specifications. It is a proposed
contract for future research runs; this editorial revision does not
supply the scientific certificates themselves.

\subsection{What ``pure thought'' means here}

Use explicitly stated mathematical models and deductions without fitting
conclusions to experimental outcomes. Synthetic examples, exact model
inputs, analytic references, numerical search, and established software
are allowed. Physical constants, electronic-structure approximations,
constitutive laws, and empirical parameter sets must be identified when
used. A model-derived result is conditional on that model.

The original ban on legacy software conflicts with its proposed reliance
on NumPy, SciPy, CVXPY, nauty, and other libraries. Reimplementing every
optimizer and algebra system is not necessary for scientific
independence. Prefer a small independently checkable certificate and a
documented trusted base. A fresh implementation can be a project
objective, but must not be confused with a mathematical proof or novelty
claim.

\subsection{Four evidence levels}

{\def\LTcaptype{none} % do not increment counter
\begin{longtable}[]{@{}
  >{\raggedright\arraybackslash}p{(\linewidth - 4\tabcolsep) * \real{0.3333}}
  >{\raggedright\arraybackslash}p{(\linewidth - 4\tabcolsep) * \real{0.3333}}
  >{\raggedright\arraybackslash}p{(\linewidth - 4\tabcolsep) * \real{0.3333}}@{}}
\toprule\noalign{}
\begin{minipage}[b]{\linewidth}\raggedright
Label
\end{minipage} & \begin{minipage}[b]{\linewidth}\raggedright
What it establishes
\end{minipage} & \begin{minipage}[b]{\linewidth}\raggedright
What it does not establish
\end{minipage} \\
\midrule\noalign{}
\endhead
\bottomrule\noalign{}
\endlastfoot
Numerical evidence & A reproducible computation at stated precision,
sample size, and discretization & Exact truth, continuum coverage, or a
universal theorem \\
Finite-model certificate & An exact/validated statement about the
specified finite model or finite relaxation & An omitted tail,
thermodynamic limit, or physical realization \\
Conditional mathematical theorem & A proof over the stated domain with
every hypothesis and error controlled & Truth outside those assumptions
or a demonstrated physical system \\
Formal proof & A checked theorem in a named proof assistant, with
imported axioms/dependencies listed & Correct physical modeling merely
by virtue of formalization \\
\end{longtable}
}

Formalization is an optional extra deliverable unless the specific
project makes it mandatory. High precision, JSON output, solver status
and small residuals do not by themselves confer any of these proof
levels.

\subsection{Claim directions}

For minimization over a larger feasible set, the optimum is a lower
bound on the original minimum. A physical feasible point supplies an
upper bound. For maximization the directions reverse. Track feasibility
of the primal/dual witnesses separately from their objective values.

Infeasibility of an inner ansatz excludes only that ansatz.
Infeasibility of an outer relaxation can exclude the original problem if
the map into the relaxation is proved and all approximations preserve
necessity. Passing necessary tests means ``not excluded''; existence
needs a constructive sufficient argument.

A smallest observed example is not a minimality theorem. A complete
catalog must define the universe, equivalence relation, boundaries, and
proof that pruning misses nothing. A finite-size numerical trend does
not prove an asymptotic limit.

\subsection{Minimum run manifest}

Before calculation, record the problem ID and revision hash, exact
inputs and conventions, quantified domain, equivalence relation, target
observable, accepted evidence level, tolerances, and resource limits.
Record dependencies and precision/rounding mode. Where randomization is
used, preserve seeds, trial counts and raw sufficient statistics.
Resource limits are planning decisions, not evidence that a mathematical
solution cannot exist.

An output bundle should include:

\begin{itemize}
\tightlist
\item
  The statement being claimed, including all assumptions and
  quantifiers.
\item
  Construction or proof object, plus an independently runnable verifier
  and its trust assumptions.
\item
  Numerical/error ledger separating roundoff, discretization,
  truncation, optimization, sampling and modeling error where relevant.
\item
  Positive reference cases and negative controls that the verifier must
  reject.
\item
  Coverage: included domain, omitted domain, unresolved regions and
  termination reason.
\end{itemize}

These are requirements for future implementations, not a promise of
polynomial-size certificates for every problem.

\subsection{Checking specific certificate types}

\textbf{LP/SDP/SOS:} verify equality constraints, objective direction,
and cone feasibility using exact arithmetic or outward-rounded
intervals. Rationalizing a numerical solution can break PSD constraints.
Correct residuals with justified norm bounds; do not silently project
and reuse the old objective. Distinguish a proof of a polynomial
approximation from a proof for the original function.

\textbf{Interval numerics:} state function enclosures, domains and
outward rounding. A small root box needs existence and, if claimed,
uniqueness. Cover the complement for a global enumeration. Include
singularities and tails rather than discarding them numerically.

\textbf{SAT/SMT:} a solver's UNSAT label is not a certificate. Preserve
a proof supported by a checker and the proof that the encoding, symmetry
reduction and blocked solutions match the mathematical universe. DRAT
applies to an appropriate Boolean CNF proof, not automatically to every
real-arithmetic or SDP claim.

\textbf{Statistics:} specify the estimand, dependence assumptions,
confidence procedure and bias. Zero observed failures does not prove
zero failure probability. MCMC error bars based on estimated
autocorrelation are not generally rigorous mixing guarantees.

\subsection{Milestones and novelty}

Baseline means a working, independently validated result on a concrete
reference problem. Strong means a broader or sharper result with its
claimed evidence level. Research extension means a possible novel result
whose feasibility and significance need review. A negative theorem can
be valuable; an unsuccessful search can also be useful if its scope is
documented, but the two are not interchangeable.

Do not promise ``publication-quality'' in a fixed number of months.
Before claiming novelty, conduct a focused current literature comparison
against the exact theorem/model/metric. The sources in this revision
check key definitions and limitations; they are not an exhaustive
novelty audit across all 30 fields.

\end{document}
